\documentclass[10pt,a4paper]{amsart}
\usepackage[margin=19mm]{geometry}
\usepackage{amsmath,amssymb,mathtools}
\usepackage[scr=rsfs,cal=euler]{mathalfa}
\usepackage{xfrac}
\usepackage[unicode,hidelinks,bookmarks=false]{hyperref}
\newcommand{\ie}{i.e.\ }
\newcommand{\cf}{cf.\ }
\newcommand{\resp}{respectively\ }
\newcommand{\Subsection}{\subsection}
\providecommand{\nobreakdash}{\nobreak\hbox{-}}
\setlength{\emergencystretch}{2em}
\multlinegap0pt
\usepackage{amssymb}
\usepackage{mathtools}
\usepackage{xcolor}
\usepackage[normalem]{ulem}
\usepackage{dsfont}
\newtheorem{assumption}{Assumption}
\newtheorem{thm}{Theorem}
\newtheorem{cor}{Corollary}
\title{Spatial Coagulation Systems with Mercer Kernels: Replicator Dynamics and Gelation}
\author{Ivan Kryven, Elena Magnanini}
\date{Source: arXiv:2606.22967v1 (CC BY 4.0)}
\begin{document}
\maketitle

\noindent\emph{Source and licence.} Spatial Coagulation Systems with Mercer Kernels: Replicator Dynamics and Gelation. Version arXiv:2606.22967v1 (CC BY 4.0), distributed by arXiv under the Creative Commons Attribution 4.0 International licence, http://creativecommons.org/licenses/by/4.0/, as recorded in the arXiv licence metadata for this exact version and shown on the abstract page https://arxiv.org/abs/2606.22967v1. Full licence notice: \texttt{licenses/arxiv-2606.22967-CC-BY-4.0.txt}.\par\smallskip

\noindent\textit{Editorial scope.} Counted unit: the article's cor cor:sharp-bounds (source0.tex lines 488--501). The article's complete original proof of it is delivered with it (source0.tex lines 1100--1201, one \texttt{proof} environment of 102 lines, which the article places away from the statement). No mathematics is written, repaired or re-derived: the statement and the proof are the article's own lines, in the article's own notation, and no retained line is substituted or re-flowed.  Retained uncounted as the source-local context the proof needs: lines 350--376; lines 454--479; lines 467--469.  Nothing was omitted from any retained span, so the package carries no omission record.  No retained line contains a citation, so the wrapper carries no bibliography and no bibliography entry.  No retained line contains a figure environment, an image-inclusion command or drawing code, so no image is delivered and the package declares no asset dependency.  Editorial formatting only, in addition: an AMS \texttt{amsart} preamble that restates the article's own declarations and macros used by the retained lines, namely the environments \texttt{assumption}, \texttt{thm}, \texttt{cor} on separate counters, together with the standard packages those lines need.  The retained environments print in this document as Assumption 1, Theorem 1, Corollary 1 in order of appearance, whereas the article prints the same items with the numbering of its own full document.  The complete native source of the article as distributed in the arXiv e-print for this exact version is bundled unchanged as \texttt{sources/203394/source0.tex}.
\begin{assumption}\label{ass:moment-justification}
There exists $T_1>0$ such that the following hold for all $t\in[0,T_1)$.

\smallskip
\noindent\hypertarget{ass:A1}{\textnormal{(A1)}}  
For every $x\in\mathcal S$, the series


\[
\sum_{y\in\mathcal S} G(x,y)\,m_2(x,t)\,m_1(y,t)
\]


converges uniformly in $t$.

\smallskip
\noindent\hypertarget{ass:A2}{\textnormal{(A2)}}  
For every $x\in\mathcal S$, the series


\[
\sum_{y\in\mathcal S} G(x,y)\,m_2(x,t)\,m_2(y,t)
\]


converges uniformly in $t$.
\end{assumption}

\begin{thm}[Gelation time]
\label{thm:riccati-general}
Assume Assumption~\ref{ass:moment-justification}.
For every $t$ such that $M_2(t)\in(0,\infty)$, define normalised local second moment
\begin{equation}\label{eq:pt}
p_t(x)=\frac{m_2(x,t)}{M_2(t)}.
\end{equation}
Then $M_2(t)$ satisfies 
\begin{equation}
\label{eq:Riccati-general}
\frac{d}{dt}M_2(t)=\alpha(t)M_2(t)^2
\end{equation}
where
\begin{equation}\label{eq:alfa}
\alpha(t)=\sum_{x,y\in\mathcal S} G(x,y)p_t(x)p_t(y).
\end{equation}
Therefore the gelation time is given by
$
t_{\mathrm{gel}}
=
\inf
\Big\{
t>0:\int_0^t\alpha(s)ds=M_2(0)^{-1}
\Big\}.
$
\end{thm}

\begin{equation}
\alpha(t)=\sum_{x,y\in\mathcal S} G(x,y)p_t(x)p_t(y).
\end{equation}

\begin{cor}[Bounds on $t_{\rm gel}$]
\label{cor:sharp-bounds} Assume
Assumption~\ref{ass:moment-justification}. Then, for all $t$ such that
$M_2(t)\in(0,\infty)$,
\begin{equation}\label{eq:alpha_bound}
\inf_{x,y\in\mathcal S} G(x,y)\le \alpha(t)\le \sup_{x,y\in\mathcal S} G(x,y),
\end{equation}
and consequently,
\begin{equation}\label{bounds_Tgel}
\big(\sup_{x,y\in\mathcal S} G(x,y)M_2(0)\big)^{-1} \le t_{\mathrm{gel}} \le \big(\inf_{x,y\in\mathcal S} G(x,y)M_2(0)\big)^{-1},
\end{equation}
with the convention that the upper bound is $+\infty$ if
$\inf_{x,y\in\mathcal S} G(x,y)=0$. 
\end{cor}

\begin{proof}[Proof of Corollary~\ref{cor:sharp-bounds}]
By Theorem~\ref{thm:riccati-general} (see Eq.~\eqref{eq:alfa}),
$
\alpha(t)=\sum_{x,y\in\mathcal S} G(x,y)p_t(x)p_t(y),
$
where $p_t$ is a probability distribution on $\mathcal S$. Hence
$p_t(x)p_t(y)\ge 0$ for all $x,y\in\mathcal S$, and
$$
\sum_{x,y\in\mathcal S} p_t(x)p_t(y)
=
\Bigl(\sum_{x\in\mathcal S} p_t(x)\Bigr)
\Bigl(\sum_{y\in\mathcal S} p_t(y)\Bigr)
=1.
$$
Therefore $\alpha(t)$ is a convex combination of the values
$\{G(x,y):x,y\in\mathcal S\}$, so that
$$
\inf_{x,y\in\mathcal S} G(x,y)
\le
\alpha(t)
\le
\sup_{x,y\in\mathcal S} G(x,y).
$$
Set
$
a:=\inf_{x,y\in\mathcal S} G(x,y)$,
$
b:=\sup_{x,y\in\mathcal S} G(x,y).
$
Then Theorem~\ref{thm:riccati-general} gives
$$
\frac{d}{dt}M_2(t)=\alpha(t)M_2(t)^2,
$$
hence
\begin{equation}\label{ineq:dM2}
aM_2(t)^2
\le
\frac{d}{dt}M_2(t)
\le
bM_2(t)^2.
\end{equation}
We first prove the lower bound on $t_{\mathrm{gel}}$. Since
$$
\frac{d}{dt}\Bigl(\frac{1}{M_2(t)}\Bigr)
=
-\frac{1}{M_2(t)^2}\frac{d}{dt}M_2(t),
$$
the inequalities in \eqref{ineq:dM2} imply
$
-b
\le
\frac{d}{dt}\Bigl(\frac{1}{M_2(t)}\Bigr)
\le
-a.
$
Integrating from $0$ to $t<t_{\mathrm{gel}}$, we obtain
$$
\frac{1}{M_2(0)}-bt
\le
\frac{1}{M_2(t)}
\le
\frac{1}{M_2(0)}-at.
$$
Now, for every $t<t_{\mathrm{gel}}$, one has $M_2(t)<\infty$, hence
$\frac1{M_2(t)}>0$. Therefore the quantity
$\frac{1}{M_2(0)}-bt$ must remain strictly below the first possible time at which
$\frac1{M_2(t)}$ can vanish. It follows that
$$
t_{\mathrm{gel}}
\ge
\frac{1}{bM_2(0)}
=
\frac{1}{\sup_{x,y\in\mathcal S} G(x,y)M_2(0)}.
$$
We next prove the upper bound. By the characterisation of the gelation time in
Theorem~\ref{thm:riccati-general},
$$
t_{\mathrm{gel}}
=
\inf
\Bigl\{
t>0:\int_0^t \alpha(s)ds=M_2(0)^{-1}
\Bigr\}.
$$
Since $\alpha(s)\ge a$ for all $s$, we have
$
\int_0^t \alpha(s)ds \geq at.
$
Hence, if $a>0$, then at time
$
t=\frac{1}{aM_2(0)}
$
the integral has already reached $M_2(0)^{-1}$, and therefore
$$
t_{\mathrm{gel}}
\le
\frac{1}{aM_2(0)}
=
\frac{1}{\inf_{x,y\in\mathcal S} G(x,y)M_2(0)}.
$$
If $a=0$, the right-hand side is understood as $+\infty$, so the estimate remains valid. This concludes the proof.
\end{proof}

\end{document}
